\documentclass[11pt]{article}
\newcommand{\ListaTitulo}{Lista 04: Correlação, assimetria e curtose}
% Preâmbulo compartilhado: MATF14, Listas de Exercícios 2026
% Usado por ListaNN.tex e GabaritoNN.tex via % Preâmbulo compartilhado: MATF14, Listas de Exercícios 2026
% Usado por ListaNN.tex e GabaritoNN.tex via % Preâmbulo compartilhado: MATF14, Listas de Exercícios 2026
% Usado por ListaNN.tex e GabaritoNN.tex via \input{preamble}

\usepackage[utf8]{inputenc}
\usepackage[T1]{fontenc}
\usepackage[brazil]{babel}
\usepackage{fullpage}
\usepackage{amsmath,amssymb}
\usepackage{enumitem}
\usepackage{xcolor}
\usepackage{fancyhdr}
\usepackage{titlesec}
\usepackage{listings}
\usepackage{hyperref}
\usepackage{booktabs}
\usepackage{graphicx}
\usepackage{tikz}

\definecolor{matf14azul}{HTML}{0D3B54}
\definecolor{matf14dourado}{HTML}{C9971E}

\titleformat{\section}{\normalfont\Large\bfseries\color{matf14azul}}{\thesection}{1em}{}
\titleformat{\subsection}{\normalfont\large\bfseries\color{matf14azul}}{\thesubsection}{1em}{}

\providecommand{\ListaTitulo}{Lista de Exercícios}

\setlength{\headheight}{14pt}
\pagestyle{fancy}
\fancyhf{}
\lhead{\small MATF14: Estatística Econômica I}
\rhead{\small \ListaTitulo}
\cfoot{\thepage}
\renewcommand{\headrulewidth}{0.4pt}

\lstset{
  basicstyle=\ttfamily\small,
  breaklines=true,
  frame=single,
  columns=fullflexible,
  backgroundcolor=\color{gray!8},
  commentstyle=\color{matf14dourado},
  keywordstyle=\color{matf14azul}\bfseries,
  language=R
}

\newcounter{questaocounter}
\newenvironment{questao}[1]{%
  \stepcounter{questaocounter}%
  \bigskip\noindent\textbf{Questão #1.}\ %
}{\par}

\newcommand{\cabecalho}[2]{%
  \begin{center}
    {\Large\bfseries\color{matf14azul} #1}\\[0.3em]
    {\large #2}
  \end{center}
  \vspace{0.5em}
  \noindent\rule{\linewidth}{0.6pt}
  \vspace{0.5em}
}

\newenvironment{solucao}{%
  \par\smallskip\noindent\textit{\color{matf14dourado!80!black}Solução.}\ %
}{\hfill$\blacksquare$\par}


\usepackage[utf8]{inputenc}
\usepackage[T1]{fontenc}
\usepackage[brazil]{babel}
\usepackage{fullpage}
\usepackage{amsmath,amssymb}
\usepackage{enumitem}
\usepackage{xcolor}
\usepackage{fancyhdr}
\usepackage{titlesec}
\usepackage{listings}
\usepackage{hyperref}
\usepackage{booktabs}
\usepackage{graphicx}
\usepackage{tikz}

\definecolor{matf14azul}{HTML}{0D3B54}
\definecolor{matf14dourado}{HTML}{C9971E}

\titleformat{\section}{\normalfont\Large\bfseries\color{matf14azul}}{\thesection}{1em}{}
\titleformat{\subsection}{\normalfont\large\bfseries\color{matf14azul}}{\thesubsection}{1em}{}

\providecommand{\ListaTitulo}{Lista de Exercícios}

\setlength{\headheight}{14pt}
\pagestyle{fancy}
\fancyhf{}
\lhead{\small MATF14: Estatística Econômica I}
\rhead{\small \ListaTitulo}
\cfoot{\thepage}
\renewcommand{\headrulewidth}{0.4pt}

\lstset{
  basicstyle=\ttfamily\small,
  breaklines=true,
  frame=single,
  columns=fullflexible,
  backgroundcolor=\color{gray!8},
  commentstyle=\color{matf14dourado},
  keywordstyle=\color{matf14azul}\bfseries,
  language=R
}

\newcounter{questaocounter}
\newenvironment{questao}[1]{%
  \stepcounter{questaocounter}%
  \bigskip\noindent\textbf{Questão #1.}\ %
}{\par}

\newcommand{\cabecalho}[2]{%
  \begin{center}
    {\Large\bfseries\color{matf14azul} #1}\\[0.3em]
    {\large #2}
  \end{center}
  \vspace{0.5em}
  \noindent\rule{\linewidth}{0.6pt}
  \vspace{0.5em}
}

\newenvironment{solucao}{%
  \par\smallskip\noindent\textit{\color{matf14dourado!80!black}Solução.}\ %
}{\hfill$\blacksquare$\par}


\usepackage[utf8]{inputenc}
\usepackage[T1]{fontenc}
\usepackage[brazil]{babel}
\usepackage{fullpage}
\usepackage{amsmath,amssymb}
\usepackage{enumitem}
\usepackage{xcolor}
\usepackage{fancyhdr}
\usepackage{titlesec}
\usepackage{listings}
\usepackage{hyperref}
\usepackage{booktabs}
\usepackage{graphicx}
\usepackage{tikz}

\definecolor{matf14azul}{HTML}{0D3B54}
\definecolor{matf14dourado}{HTML}{C9971E}

\titleformat{\section}{\normalfont\Large\bfseries\color{matf14azul}}{\thesection}{1em}{}
\titleformat{\subsection}{\normalfont\large\bfseries\color{matf14azul}}{\thesubsection}{1em}{}

\providecommand{\ListaTitulo}{Lista de Exercícios}

\setlength{\headheight}{14pt}
\pagestyle{fancy}
\fancyhf{}
\lhead{\small MATF14: Estatística Econômica I}
\rhead{\small \ListaTitulo}
\cfoot{\thepage}
\renewcommand{\headrulewidth}{0.4pt}

\lstset{
  basicstyle=\ttfamily\small,
  breaklines=true,
  frame=single,
  columns=fullflexible,
  backgroundcolor=\color{gray!8},
  commentstyle=\color{matf14dourado},
  keywordstyle=\color{matf14azul}\bfseries,
  language=R
}

\newcounter{questaocounter}
\newenvironment{questao}[1]{%
  \stepcounter{questaocounter}%
  \bigskip\noindent\textbf{Questão #1.}\ %
}{\par}

\newcommand{\cabecalho}[2]{%
  \begin{center}
    {\Large\bfseries\color{matf14azul} #1}\\[0.3em]
    {\large #2}
  \end{center}
  \vspace{0.5em}
  \noindent\rule{\linewidth}{0.6pt}
  \vspace{0.5em}
}

\newenvironment{solucao}{%
  \par\smallskip\noindent\textit{\color{matf14dourado!80!black}Solução.}\ %
}{\hfill$\blacksquare$\par}


\begin{document}

\cabecalho{Lista de Exercícios 04}{Coeficiente de correlação, assimetria e curtose: fecha a Unidade 1 (Cap. 1, Aula 08)}

\noindent \textit{Justifique todas as respostas. As resoluções são discutidas em sala e nos horários de atendimento; não são publicadas neste material.}

\begin{questao}{1} \textbf{(Interpretação de $r$: gastos com propaganda e vendas)}

Uma rede de varejo mediu, em 8 lojas, o gasto mensal com propaganda local (em R\$ mil) e o
faturamento do mês seguinte (em R\$ mil), obtendo $r = 0{,}86$.

\begin{enumerate}[label=(\alph*)]
  \item Interprete o valor $r=0{,}86$ em palavras, sem usar jargão técnico.
  \item O gerente de marketing conclui: "isso prova que gastar mais em propaganda causa mais
    vendas". Aponte uma explicação alternativa (não causal) plausível para essa correlação.
  \item Proponha um desenho de coleta de dados (não precisa ser sofisticado) que ajudaria a
    distinguir a hipótese causal da alternativa que você deu em (b).
\end{enumerate}
\end{questao}

\begin{questao}{2} \textbf{(Correlação quase nula não é ausência de relação)}

Um analista calcula a correlação entre a taxa de juros básica (Selic) e o volume de vendas de
imóveis ao longo de 3 anos e obtém $r \approx 0{,}03$, praticamente zero. Ele conclui que "juros
não afetam o mercado imobiliário".

\begin{enumerate}[label=(\alph*)]
  \item $r$ mede associação de que tipo especificamente? Isso limita a conclusão do analista?
  \item Descreva (em palavras, sem fórmula) um padrão nos dados, envolvendo juros e vendas de
    imóveis, que produziria $r\approx 0$ mesmo havendo uma relação forte e sistemática entre as
    duas variáveis. (Dica: pense em uma relação que não seja uma linha reta.)
\end{enumerate}
\end{questao}

\begin{questao}{3} \textbf{(Assimetria: comparando duas distribuições salariais)}

Duas empresas, X e Y, têm a mesma média salarial (R\$ 5.000) e o mesmo desvio padrão (R\$ 1.500).
A empresa X tem assimetria $\approx 0$; a empresa Y tem assimetria $\approx +1{,}8$.

\begin{enumerate}[label=(\alph*)]
  \item Descreva, em palavras, a diferença esperada no formato do histograma de salários entre X e
    Y, mesmo com média e desvio padrão idênticos.
  \item Em qual das duas empresas a mediana salarial provavelmente está mais abaixo da média?
    Justifique.
  \item Um sindicato de trabalhadores da empresa Y, ao negociar reajuste geral em R\$ (valor fixo
    igual para todos), quer saber se isso reduz ou aumenta a assimetria da distribuição salarial.
    Justifique sua resposta.
\end{enumerate}
\end{questao}

\begin{questao}{4} \textbf{(Curtose e risco financeiro)}

Os retornos diários de duas ações, P e Q, têm a mesma média e o mesmo desvio padrão histórico. A
ação P tem curtose $\approx 3$ (próxima da Normal); a ação Q tem curtose $\approx 7$.

\begin{enumerate}[label=(\alph*)]
  \item Um gestor de risco que assume (erroneamente) que os retornos de Q seguem uma distribuição
    Normal está superestimando ou subestimando a frequência de quedas bruscas ("caudas gordas")?
  \item Explique, em uma frase, por que curtose alta é uma preocupação maior para quem administra
    risco de uma carteira do que para quem só olha a média de retorno.
\end{enumerate}
\end{questao}

\begin{questao}{5} \textbf{(Dedução: invariância de $r$ à mudança de unidade)}

O coeficiente de correlação entre $X$ e $Y$ é
$$
r_{XY} = \frac{\sum_{i=1}^n (x_i-\bar x)(y_i-\bar y)}{\sqrt{\sum_i (x_i-\bar x)^2}\sqrt{\sum_i (y_i-\bar y)^2}}.
$$
Suponha que $X$ seja medida em uma nova unidade, $X' = aX + b$, com $a>0$ (por exemplo,
$X$ em dólares e $X'=5X$ em reais, a uma cotação fixa de 5).

\begin{enumerate}[label=(\alph*)]
  \item Mostre que $\overline{X'} = a\bar x + b$ e que $x_i' - \overline{X'} = a(x_i - \bar x)$.
  \item Usando o item (a), mostre que $r_{X'Y} = r_{XY}$, ou seja, o coeficiente de correlação
    \textbf{não muda} quando se troca a unidade de medida de uma das variáveis (desde que a
    transformação seja linear com $a>0$). Isso justifica por que $r$ entre "gasto em dólares" e
    "vendas" é o mesmo que entre "gasto em reais" e "vendas", à mesma cotação.
\end{enumerate}
\end{questao}

\bigskip
\noindent\textit{A questão a seguir não tem resposta única e não é cobrada em prova no mesmo
formato, fecha a Unidade 1 e é para discussão em grupo ou por escrito, antes da próxima aula.}

\begin{questao}{6 (discussão: síntese da Unidade 1)} \textbf{(Sem resposta única)}

Escolha uma das duas bases do curso (\texttt{empresas} ou \texttt{dados\_prefeitos}) e, sem usar
nenhuma fórmula nova, escreva um parágrafo de meia página descrevendo a variável de sua escolha
usando \emph{todo} o vocabulário da Unidade 1: tipo de variável, tabela ou gráfico apropriado,
tendência central, dispersão, e, se fizer sentido para a variável escolhida, assimetria e
correlação com outra variável da base. O objetivo é praticar a integração dos conceitos, não
calcular algo novo.
\end{questao}


\bigskip
\section*{Exercícios complementares}

\begin{enumerate}[label=\textbf{C\arabic*.}, leftmargin=*]
\item A tabela traz o PIB per capita (R\$) e o IDH das capitais do Nordeste.
\begin{center}
\begin{tabular}{l|ccccccccc}
\toprule
Capital & REC & AJU & SLZ & NAT & JPA & SSA & FOR & THE & MCZ \\
\midrule
PIB per capita & 33.094 & 27.139 & 36.827 & 28.409 & 21.662 & 28.483 & 24.296 & 22.279 & 24.322 \\
IDH & 0,772 & 0,770 & 0,768 & 0,763 & 0,763 & 0,759 & 0,754 & 0,751 & 0,721 \\
\bottomrule
\end{tabular}
\end{center}
\begin{enumerate}[label=(\alph*)]
  \item Faça o gráfico de dispersão.
  \item Calcule o coeficiente de correlação mostrando a tabela com os valores padronizados.
  \item Interprete o resultado. Ele permite dizer que aumentar o PIB per capita aumenta o IDH?
\end{enumerate}
\item Com os dados de 10 casais (salário anual em salários mínimos) de um bairro:
\begin{center}
\begin{tabular}{l|cccccccccc}
\toprule
Casal & 1 & 2 & 3 & 4 & 5 & 6 & 7 & 8 & 9 & 10 \\
\midrule
Homem & 10 & 10 & 10 & 15 & 15 & 15 & 15 & 20 & 20 & 20 \\
Mulher & 5 & 10 & 10 & 5 & 10 & 10 & 15 & 10 & 10 & 15 \\
\bottomrule
\end{tabular}
\end{center}
\begin{enumerate}[label=(\alph*)]
  \item Calcule, separadamente, média, moda, mediana e variância dos salários de homens e de mulheres.
  \item Calcule a correlação entre o salário do homem e o da mulher em cada casal. Interprete.
  \item Calcule a assimetria dos salários das mulheres. A distribuição é simétrica?
\end{enumerate}
\item Para os dados $X = -2, -1, 0, 1, 2$ e $Y = X^2$, calcule $cor(X, Y)$. O resultado significa que $X$ e $Y$ não estão associados? Explique.

\end{enumerate}

\bigskip
\needspace{12\baselineskip}
\section*{Aprofundamento: contas nacionais e dados reais}
\noindent\textit{Exercícios com dados oficiais, ligados à seção de contas nacionais do Capítulo 1 do livro.}

\begin{enumerate}[label=\textbf{A\arabic*.}, leftmargin=*]
\item \textbf{(Poupança e investimento.)} Taxa de poupança bruta e taxa de investimento (formação bruta de capital), em \% do PIB
(IBGE, Contas Econômicas Trimestrais, somas anuais):
\begin{center}
\begin{tabular}{lccccccc}
\toprule
Ano & 2019 & 2020 & 2021 & 2022 & 2023 & 2024 & 2025 \\
\midrule
Poupança ($x$) & 12,2 & 14,8 & 17,1 & 15,8 & 15,0 & 14,1 & 14,4 \\
Investimento ($y$) & 15,5 & 16,1 & 19,5 & 18,1 & 15,8 & 17,0 & 17,1 \\
\bottomrule
\end{tabular}
\end{center}
\begin{enumerate}[label=(\alph*)]
  \item Faça o diagrama de dispersão e calcule o coeficiente de correlação de Pearson.
  \item Em todos os anos o investimento superou a poupança. Pela identidade das contas econômicas integradas, como essa diferença foi financiada?
  \item Correlação alta entre poupança e investimento prova que ``a poupança causa o investimento''? Dê uma explicação alternativa.
\end{enumerate}

\item \textbf{(Assimetria e curtose do crescimento do PIB.)} Para as 121 taxas de variação do PIB em relação ao trimestre anterior (2º trimestre de 1996 ao
2º trimestre de 2026), obtêm-se média 0,58\%, assimetria $-1{,}18$ e curtose $16{,}0$. Retirando os quatro trimestres de 2020 ($-2{,}3$; $-8{,}9$; $7{,}9$; $3{,}8$),
a assimetria passa a $-0{,}76$ e a curtose a $5{,}0$.
\begin{enumerate}[label=(\alph*)]
  \item Interprete o sinal da assimetria. O que ele diz sobre recessões e expansões?
  \item Compare a curtose com o valor de referência 3. Por que quatro observações mudam tanto a curtose?
  \item Para descrever o crescimento ``típico'' da economia, você usaria média e desvio padrão ou mediana e distância interquartílica? Justifique.
\end{enumerate}

\end{enumerate}

\end{document}
